\par
\SourceLeaf{378}%
We think we ought to add a few remarks. We find that in every measurement the index was always brought back to 0 at the start of each series. We cannot approve this procedure, since it means that however many more series are taken, the same constant error due to the error of graduation must always remain in the result. --- The condition of the air is stated for each of the Piedmontese officers' measurements. Of 414 series of measurements, we count 320 in which calm and 94 in which wind is reported. Such a favourable proportion would hardly have been expected at stations so high up (the highest is 3534 metres above sea level).

The second volume contains the astronomers' work in ten sections. The first two sections contain determinations of differences in longitude by gunpowder signals, relating to the measurement of the arc of longitude. The first trials of this kind were made in September 1821. The signals were given on Rocca Melone and observed at the Milan observatory, 170,000 metres away, and on the nearby Mont Cenis. To determine the time at the latter site, a small observatory had been built in the hospice garden and a transit telescope by \textsc{Fortin} set up in it. The instrument, however, appears to have been of less than outstanding quality, as is expressly remarked in connection with the pivots, an essential part of it. Rocca Melone could not be seen from here. To see the signals, the observers had to go to a spot some distance away, to which they carried the time with an \textsc{Earnshaw} chronometer. The transit telescope observations, too, were usually recorded against this chronometer, by an assistant at a sign from the observer himself. All these circumstances, to be sure, combine to diminish confidence in the. accuracy of the final result, even though the three partial results from the three days of observation agree very well. In addition, time was determined here from stars and in Milan from transits of the Sun. Finally, it seems that the perpendicularity of the optical axis of \textsc{Fortin}'s transit telescope to its axis of rotation was left entirely uncorrected. At any rate, this important matter is never mentioned in connection with these observations.

In similar operations in 1822,
